Applying the periodic doubling argument of Section~\ref{sm:sec:trotter} to a
new Hamiltonian requires two inputs:
\begin{enumerate}
    \item write $H_L=\sum_j h_{j,j+1}+\sum_j o_j$ in the real
          bond-reflection-positive form $-h=\sum_\alpha C_\alpha\otimes FC_\alpha F$
          with $FoF=o$ on a chosen unit cell, so that the two-cell Trotter
          transfer matrix $\smT=\smF^{\mathsf T}\smF\succeq0$ of
          Eq.~\eqref{sm:eq:positive} is positive semidefinite at every even cell
          count $L$;
    \item bound $1-\Ps(n_0,\beta_0)$ and $1-\Pt(2n_0,\beta_0)$ at one
          finite scale by rationals $(p_0,q_0)$ inside the convergence region of
          the doubling recurrence.
\end{enumerate}

Reflection positivity carries no gap information (gapless and
spontaneously dimerized chains are reflection positive too;
Proposition~\ref{models:cor:obstruct}), so the seed computation must be
matched to the correlation length $\xi$ and gap $\Delta$ of the model:
\begin{center}
    \renewcommand{\arraystretch}{1.3}
    \setlength{\tabcolsep}{4pt}
    \begin{tabular}{@{}llp{4.1cm}p{5.3cm}@{}}
        \toprule
        Regime                           & Finite inputs          & Seed strategy                           & Models in Table~\ref{tab:models}                                     \\
        \midrule
        Short $\xi$, small $d^{2n_0}$    & $2n_0\le16$ spins      & Symmetry blocks + Weyl bound            & Staggered, dimerized, ladders, Ising (Sec.~\ref{models:sec:spectra}) \\
        Moderate $\xi$, large $d^{4n_0}$ & $6,12$ spins           & Block spectra + Eq.~\eqref{sm:eq:torus} & Single-ion anisotropy, $D=4$ (Sec.~\ref{models:sec:spectra})         \\
        Weak off-diagonal term           & $6,12$ spins           & Trial residuals + diagonal count        & Mixed-field Ising, kernel seed (Sec.~\ref{models:sec:ising})         \\
        Long $\xi$ ($n_0\ge36$)          & $4$--$12$ traces + MPS & Twisted-ring moments + integer MPS      & Heisenberg, $S=1$ (Sec.~\ref{sm:sec:seeds})                          \\
        \bottomrule
    \end{tabular}
\end{center}

\subsection{Reflection-positive spin chains, unit cells, and obstructions}
\label{models:sec:rp}

To apply the finite-Trotter Gram factorization of
Section~\ref{sm:sec:transfer} to a periodic ring of $L\ge2$ unit cells,
\begin{equation}
    H_L=\sum_{j=1}^L h_{j,j+1}+\sum_{j=1}^L o_j,
    \qquad j+L\equiv j,
    \label{models:eq:ring}
\end{equation}
we require a real local cell basis, a real symmetric involution $F$
($F^2=\smI$), a real symmetric on-site matrix $o$ satisfying $FoF=o$,
and real symmetric or antisymmetric bond matrices $C_\alpha$ such that
\begin{equation}
    -h=\sum_\alpha C_\alpha\otimes FC_\alpha F
    =\sum_\alpha\;\scopebond\,.
    \label{models:eq:rp}
\end{equation}
Hermiticity of $h$ alone does not supply this form: the factors must be
real and $o$ even under $F$, so basis and cell are chosen model by model.

\subsubsection{Bilinear, staggered, and biquadratic spin chains}
\label{models:sec:catalogue}

For a bilinear bond, a factor $v\cdot\bm S$ and its mirror image must
multiply back to $-\bm S^{\mathsf T}J\bm S$, so the mirror has to act on
the spin components by a $3\times3$ matrix $O$. Only maps realized by an
antiunitary symmetry qualify, which restricts the sign pattern of $J$.
Consider first a single-site spin-$S$ chain
($F=\smI$) with a real symmetric exchange matrix $J$ and a uniform
field $\bm h$,
\begin{equation}
    H_L=\sum_{j=1}^L\bm S_j^{\mathsf T}J\bm S_{j+1}
    +\sum_{j=1}^L\bm h\cdot\bm S_j.
    \label{models:eq:bilinear}
\end{equation}
In the standard $S^z$ eigenbasis, $S^x$ and $S^z$ are real and $S^y$ is
purely imaginary, so complex conjugation $K$ acts on the spin vector by
$\operatorname{diag}(1,-1,1)$. Composing $K$ with a $\pi$ rotation $U_R$
about a principal axis gives an antiunitary reflection $\Theta=U_RK$ that
acts on the spin vector by a symmetric orthogonal matrix
$O=R\operatorname{diag}(1,-1,1)$ with $\det O=-1$; the rotation about $y$
gives time reversal $T=e^{i\pi S^y}K$ with $O=-I$.

\begin{proposition}[Bilinear spin chains]
    \label{models:prop:xyz}
    If there exists a symmetric orthogonal matrix $O\in\mathrm O(3)$ with
    $\det O=-1$ such that $-JO\succeq0$ and $O\bm h=\bm h$, then
    Eq.~\eqref{models:eq:bilinear} admits a reflection-positive bond
    factorization for every spin $S$. For diagonal exchange
    $J=\operatorname{diag}(J_x,J_y,J_z)$ with nonzero entries and $\bm h=0$,
    such an $O$ exists if and only if an even number of $J_x,J_y,J_z$ are
    negative.
\end{proposition}

\begin{proof}
    Factor the positive semidefinite matrix $-JO=\sum_\alpha v_\alpha v_\alpha^{\mathsf T}$
    with real vectors $v_\alpha\in\mathbb R^3$ and set
    $C_\alpha=v_\alpha\cdot\bm S$. Under $\Theta=U_RT$,
    $\Theta(C_\alpha)=(Ov_\alpha)\cdot\bm S$, so
    $\sum_\alpha C_\alpha\otimes\Theta(C_\alpha)
        =\bm S^{\mathsf T}(-JO)O\bm S=-\bm S^{\mathsf T}J\bm S$, while
    $\Theta(\bm h\cdot\bm S)=(O\bm h)\cdot\bm S=\bm h\cdot\bm S$. When $J$
    is diagonal with nonzero entries, $-JO\succeq0$ forces
    $O=\operatorname{sign}(-J)$, whose determinant $(-1)^{n_-}$ is $-1$ if
    and only if the number $n_-$ of ferromagnetic ($J_\alpha<0$) components
    is even.
\end{proof}

Proposition~\ref{models:prop:xyz} sorts the standard single-site spin
chains by the sign parity $n_-$ of their exchange couplings; a uniform
field fails because $O=-\smI$ reverses it:
\begin{center}
    \setlength{\tabcolsep}{5pt}
    \begin{tabular}{@{}llll@{}}
        \toprule
        Chain                                         & $(J_x,J_y,J_z)$, $\bm h$ & Mirror $O$                                  & Reflection positive \\
        \midrule
        Heisenberg antiferromagnet, any $S$           & $(1,1,1)$, $0$           & $-\smI$ ($n_-=0$)                           & yes                 \\
        XXZ, $\Delta\ge0$                             & $(1,1,\Delta)$, $0$      & $-\smI$                                     & yes                 \\
        XY, either sign                               & $(J_x,J_y,0)$, $0$       & $J_z=0$ frees $O_{zz}$                      & yes                 \\
        Transverse- or mixed-field Ising              & $(0,0,-1)$, $(g,0,h)$    & $\operatorname{diag}(1,-1,1)$ fixes $\bm h$ & yes                 \\
        Heisenberg ferromagnet                        & $(-1,-1,-1)$, $0$        & $n_-=3$                                     & no                  \\
        XXZ, $\Delta<0$                               & $(1,1,\Delta)$, $0$      & $n_-=1$                                     & no                  \\
        Heisenberg antiferromagnet in a uniform field & $(1,1,1)$, $\bm h\ne0$   & $-\smI$ reverses $\bm h$                    & no                  \\
        \bottomrule
    \end{tabular}
\end{center}

Although a uniform field breaks reflection positivity of the Heisenberg
antiferromagnet, a \emph{staggered} field restores it on every even
ring: a $\pi$ rotation about the $y$ axis on every odd site makes the
field uniform,
\begin{equation}
    \staggerrotation\qquad
    -h_{j,j+1}=\sum_\alpha\;\rpbondpic{P_\alpha}{P_\alpha}\,,
    \label{models:eq:stagger-rotation}
\end{equation}
and the bond sign parity stays even, $n_-=2$, because the remaining sign
sits in the real antisymmetric factor $P_y=iS^y$. The same rotation
underlies the staggered-field chain in Table~\ref{tab:models} and, on
open chains, relates opposite endpoint fields \fieldopp\ to equal
endpoint fields \fieldsame\ (Section~\ref{sm:sec:obc}).

For $S=1$, adding a biquadratic interaction gives the family
\begin{equation}
    H_L(\vartheta)=\sum_{j=1}^L\Bigl[
        \cos\vartheta\,\bm S_j\cdot\bm S_{j+1}
        +\sin\vartheta\,(\bm S_j\cdot\bm S_{j+1})^2\Bigr].
    \label{models:eq:blbq}
\end{equation}
Proposition~\ref{models:prop:xyz} does not reach the biquadratic term,
which mixes products $S^\alpha S^\beta$ across the bond; factoring it
needs one-site real operators without cross terms across the mirror. In
the Cartesian basis $\ket{x},\ket{y},\ket{z}$ of
Eq.~\eqref{sm:eq:generators}, the spin generators $G_\alpha=iS^\alpha$
are real antisymmetric and the products $Q_{\alpha\beta}=S^\alpha S^\beta$
are real matrices whose symmetric part
$Q_{\alpha\beta}^+=\tfrac12(Q_{\alpha\beta}+Q_{\beta\alpha})$ and
antisymmetric part $\tfrac12\epsilon_{\alpha\beta\gamma}G_\gamma$ are
orthogonal under the trace. Expanding $(\bm S_1\cdot\bm S_2)^2$ into
$G_\alpha$ and $Q_{\alpha\beta}^+$ diagonalizes the bond across the
mirror:
\begin{equation}
    -h(\vartheta)=\bigl(\cos\vartheta-\tfrac12\sin\vartheta\bigr)
    \sum_\alpha\;\rpbondpic{G_\alpha}{G_\alpha}
    \;-\sin\vartheta\sum_{\alpha,\beta}\;
    \rpbondpic{Q_{\alpha\beta}^+}{Q_{\alpha\beta}^+}\,.
    \label{models:eq:blbq-decomp}
\end{equation}

\begin{proposition}[Reflection positivity of the bilinear--biquadratic chain]
    \label{models:prop:blbq}
    In the Cartesian basis, the bond~\eqref{models:eq:blbq-decomp} is
    reflection positive if and only if both coefficients are nonnegative,
    that is,
    \begin{equation}
        -\pi+\arctan2\le\vartheta\le0.
        \label{models:eq:blbq-window}
    \end{equation}
\end{proposition}

\begin{proofdiagram}
    \blbqcircle
\end{proofdiagram}
\noindent Reflection positivity of the whole family thus reduces to the
signs of two coefficients; the frustration-free AKLT point
($\tan\vartheta=1/3>0$) lies outside the arc because its positive
biquadratic term has the wrong sign in Eq.~\eqref{models:eq:blbq-decomp}.
The doubling argument applies only on the Haldane sub-arc $(-\pi/4,0]$: the
spontaneously dimerized phase $\vartheta\in[-\pi+\arctan2,-\pi/4]$ and
the gapless spin-$1/2$ Heisenberg chain are reflection positive but
have no seed in the convergence region, for reasons no larger ring or lower
temperature removes:

\begin{proposition}[Obstructions to a convergent seed]
    \label{models:cor:obstruct}
    For any reflection-positive ring Hamiltonian~\eqref{models:eq:ring}, no
    scale $(n_0,\beta_0)$ has purities $\Ps(n_0,\beta_0)$ and
    $\Pt(2n_0,\beta_0)$ inside the convergence region of the doubling
    recurrence if either
    \begin{enumerate}
        \item[(i)] $H_{2n_0}$ has a $k$-fold degenerate ground level with
              $k\ge2$: the ordered Gibbs weights satisfy $w_0=w_1\le1/k\le1/2$
              (\specbars{.49,.49,.02}), so $\Pt(2n_0,\beta_0)=\sum_m w_m^2\le w_0\le1/2$;
              or
        \item[(ii)] $\liminf_{L\to\infty,\;L\text{ even}}\Delta_L=0$: a
              seed in the convergence region would give a uniform positive lower bound on $\Delta_L$
              for all even $L\ge n_0$.
    \end{enumerate}
\end{proposition}
\noindent Reflection positivity thus supplies the squaring map, and only
a finite seed inside the convergence region separates gapped chains with
a unique ground state from critical or symmetry-broken ones.

\subsubsection{Multi-spin unit cells, dimer reflections, and explicit decompositions}
\label{models:sec:explicit-rp}

When interactions alternate from bond to bond or couple two parallel
legs, we block sites into a multi-spin unit cell. A single-site mirror
then fails, because the inter-cell bond couples the last spin of one
cell to the first spin of the next. The cell reflection $F$ repairs
this: across an inter-cell mirror cut, a reflection of the ring
reverses the left-to-right order of the sites inside each cell:
\begin{equation}
    -h_{j,j+1}=\sum_\alpha\;\cellbondform\,.
    \label{models:eq:cell-mirror}
\end{equation}
In a two-spin dimer cell $(a,b)$, $F$ is the two-spin swap
$F(x\otimes y)=y\otimes x$%
~\leanmark{https://sirui-lu.com/CertifiedGaps/docs/CertifiedGaps/Local/CellSwap.html\#CertifiedGaps.cellSwap_mul_cellSwap},
which maps $C_\alpha=\smI\otimes G_\alpha$ on spin $b$ to
$FC_\alpha F=G_\alpha\otimes\smI$ on spin $a$ of the next cell as drawn, and the intra-cell bond is a
swap-invariant on-site matrix $FoF=o$. In a two-leg ladder the cell is
a rung and the mirror keeps each leg on its own side,
\begin{equation}
    -h_{j,j+1}=\sum_\alpha\;\rungcellform\,,
    \label{models:eq:rung-mirror}
\end{equation}
so the rung interaction is even under the reflection and belongs to
$o$. These two reflections, the dimer swap and the identity on a rung,
are all that Table~\ref{models:tab:rp} needs.

The two-cell Gram form makes $Z_n(\beta)$ a sum of nonnegative powers
only at even $n$, so an odd seed length needs positivity of the one-cell
transfer operator itself. A dimer cell whose intra-cell bond is
reflection positive supplies it: cutting the cell across its internal
bond factors the \emph{one-cell} spatial transfer operator into a
half-cell product with its transpose,
\begin{equation}
    X_\beta=\oddcellpos=KK^{\mathsf T}\succeq0.
    \label{models:eq:odd-cell}
\end{equation}
Every eigenvalue of $KK^{\mathsf T}$ is nonnegative, so
$Z_n(\beta)=\Tr X_\beta^n=\sum_i\lambda_i(\beta)^n$ is a sum of
nonnegative powers at \emph{odd} cell counts $n_0$ too.
In the dimerized spin-$1$ chain at $\delta=7/10$,
Eq.~\eqref{models:eq:odd-cell} allows an odd seed of $n_0=3$ dimer cells
($6$ and $12$ spins), avoiding the $16$-spin ($3^{16}\simeq4.3\times10^7$)
ring that an even seed $n_0=4$ would require.

Proposition~\ref{models:prop:xyz} guarantees a factorization, but the
exact checkers behind Table~\ref{models:tab:rp} need its factors as
explicit real matrices, and $S^y$ is imaginary in the $S^z$ basis.
For $S=1/2$, we write $P_x=S^x=\tfrac12\sigma^x$,
$P_y=iS^y=\tfrac12(i\sigma^y)$, and $P_z=S^z=\tfrac12\sigma^z$, which
are real matrices satisfying $P_x^{\mathsf T}=P_x$,
$P_y^{\mathsf T}=-P_y$, and $P_z^{\mathsf T}=P_z$. Rotating one site of
an antiferromagnetic bond by $u_y=e^{i\pi S^y}=2P_y$ maps
$(S^x,S^y,S^z)\mapsto(-S^x,S^y,-S^z)$ and converts $-\bm S_1\cdot\bm S_2$
into the real reflection-positive form%
~\leanmark{https://sirui-lu.com/CertifiedGaps/docs/CertifiedGaps/Local/SpinHalfRotation.html\#CertifiedGaps.SpinHalf.rot_conj_spinHalf}
\begin{equation}
    -u_{y,2}(\bm S_1\cdot\bm S_2)u_{y,2}^\dagger
    =\sum_{\alpha=x,y,z}P_\alpha\otimes P_\alpha
    =\sum_{\alpha=x,y,z}\;\rpbondpic{P_\alpha}{P_\alpha}\,.
    \label{models:eq:spinhalf-bond}
\end{equation}
Applying $u_y$ on every odd site of an even spin-$1/2$ chain (or on the
second spin $b$ of every dimer cell, or on the checkerboard sublattice
of a two-leg ladder for even $L$) is a unitary transformation that
preserves the spectrum%
~\leanmark{https://sirui-lu.com/CertifiedGaps/docs/CertifiedGaps/Ring/Conjugation.html\#CertifiedGaps.RPRing.eigenvalues_eq_spectrum}
and puts every antiferromagnetic bond into the
form~\eqref{models:eq:spinhalf-bond}. For $S=1$, we work directly in the
Cartesian basis $\ket{x},\ket{y},\ket{z}$ of
Eq.~\eqref{sm:eq:generators}, where $G_\alpha=iS^\alpha$ is real
antisymmetric, $-\bm S_1\cdot\bm S_2=\sum_\alpha G_\alpha\otimes G_\alpha$,
and $(S^z)^2=\operatorname{diag}(1,1,0)$ is real diagonal without any
sublattice rotation. Table~\ref{models:tab:rp} lists the resulting cell,
reflection $F$, bond operators $C_\alpha$, and on-site matrix $o$ for all
seven additional models in Table~\ref{tab:models}.

\begin{table*}[t]
    \caption{Explicit reflection-positive cell decompositions
        $-h=\sum_\alpha C_\alpha\otimes FC_\alpha F$ and $FoF=o$ for the seven
        additional models in Table~\ref{tab:models}, with links to their
        Lean~4 reflection-positivity proofs. Here $F$ is the cell
        reflection involution, $u_y=e^{i\pi S^y}$ is a $\pi$ spin rotation about
        $y$, $P_\alpha\in\{S^x,iS^y,S^z\}$ are the real spin-$1/2$ operators of
        Eq.~\eqref{models:eq:spinhalf-bond}, $G_\alpha=iS^\alpha$ are the real
        antisymmetric spin-$1$ generators in the Cartesian basis, and
        $S_{ab}(x\otimes y)=y\otimes x$ swaps the two spins of a dimer cell.}
    \label{models:tab:rp}
    \setlength{\tabcolsep}{5pt}
    \footnotesize
    \begin{tabular}{@{}lllll@{}}
        \toprule
        Model                                                                                                                                                                   & Unit cell ($F$)           & Local basis / rotation & Bond operators $C_\alpha$                    & On-site cell matrix $o$                                               \\
        \midrule
        Staggered field, $S=\tfrac12$~\leanmark{https://sirui-lu.com/CertifiedGaps/docs/CertifiedGaps/Models/StaggeredHeisenberg.html\#CertifiedGaps.StaggeredHeisenberg.model} & $1$ spin ($\smI$)         & $u_y$ on odd sites     & $P_\alpha$ ($\alpha=x,y,z$)                  & $-h_{\mathrm{st}}P_z$                                                 \\
        Dimerized, $S=\tfrac12$~\leanmark{https://sirui-lu.com/CertifiedGaps/docs/CertifiedGaps/Models/DimerizedHeisenberg.html\#CertifiedGaps.DimerizedHeisenberg.model}       & $2$-spin dimer ($S_{ab}$) & $u_y$ on spin $b$      & $\sqrt{1-\delta}\,(\smI\otimes P_\alpha)$    & $-(1+\delta)\sum_\alpha P_\alpha\otimes P_\alpha$                     \\
        Dimerized, $S=1$~\leanmark{https://sirui-lu.com/CertifiedGaps/docs/CertifiedGaps/Models/SpinOneDimerized.html\#CertifiedGaps.SpinOneDimerized.model}                    & $2$-spin dimer ($S_{ab}$) & Cartesian basis        & $\sqrt{1+\delta}\,(\smI\otimes G_\alpha)$    & $-(1-\delta)\sum_\alpha G_\alpha\otimes G_\alpha$                     \\
        Single-ion anisotropy, $S=1$~\leanmark{https://sirui-lu.com/CertifiedGaps/docs/CertifiedGaps/Models/SpinOneLargeD.html\#CertifiedGaps.SpinOneLargeD.model}              & $1$ spin ($\smI$)         & Cartesian basis        & $G_\alpha$ ($\alpha=x,y,z$)                  & $D\operatorname{diag}(1,1,0)$                                         \\
        Heisenberg ladder, $S=\tfrac12$~\leanmark{https://sirui-lu.com/CertifiedGaps/docs/CertifiedGaps/Models/HeisenbergLadder.html\#CertifiedGaps.HeisenbergLadder.model}     & $2$-spin rung ($\smI$)    & checkerboard $u_y$     & $P_\alpha\otimes\smI,\;\smI\otimes P_\alpha$ & $-J_\perp\sum_\alpha P_\alpha\otimes P_\alpha$                        \\
        Ising ladder~\leanmark{https://sirui-lu.com/CertifiedGaps/docs/CertifiedGaps/Models/TwoLegIsing.html\#CertifiedGaps.TwoLegIsing.model}                                  & $2$-spin rung ($\smI$)    & $\sigma^z$ basis       & $\sigma^z\otimes\smI,\;\smI\otimes\sigma^z$  & $-\sigma^z\otimes\sigma^z-g(\sigma^x\otimes\smI+\smI\otimes\sigma^x)$ \\
        Mixed-field Ising~\leanmark{https://sirui-lu.com/CertifiedGaps/docs/CertifiedGaps/Models/IsingLongitudinal.html\#CertifiedGaps.IsingLongitudinal.model}                 & $1$ spin ($\smI$)         & $\sigma^z$ basis       & $\sigma^z$                                   & $-g\sigma^x-h\sigma^z$                                                \\
        \bottomrule
    \end{tabular}
\end{table*}

\subsection{Seed bounds from verified finite-ring block spectra}
\label{models:sec:spectra}

Once a model is in the reflection-positive form~\eqref{models:eq:rp},
the doubling recurrence of Section~\ref{sm:sec:trotter} applies. For
the first six models in
Table~\ref{tab:models}, the correlation length is short enough that the
seed scale $n_0$ lies between $2$ and $12$ unit cells. Even there, the
upper bounds $(p_0,q_0)$ need two-sided enclosures of three partition
functions, and a dense eigensolve returns only floating-point numbers
whose accuracy nothing certifies. The remedy is a posteriori: take the
floating-point eigenbasis as a candidate and certify it by residual
norms whose own rounding error is bounded.

To enclose the complete spectrum of $H_L$ at $L\in\{n_0,2n_0\}$ (up to
$16$ spin-$1/2$ or $12$ spin-$1$ sites), we block-diagonalize $H_L$ by
translation, lattice reflection, and conserved spin or parity quantum
numbers into self-adjoint blocks $B_q$ whose exact integer multiplicities
$m_q$ sum to $d^L$, with every entry enclosed in rational or interval
arithmetic. From a floating-point eigenbasis $V$ and candidate
eigenvalues $\tilde\lambda_i$ of $B_q$, we bound the Frobenius norms of
the residual $R=B_qV-V\operatorname{diag}(\tilde\lambda)$ and orthogonality
error $G=V^{\mathsf T}V-\smI$ (including floating-point rounding
error~\cite{Rump2010Verification}); Ostrowski's and Weyl's theorems then
enclose the sorted exact eigenvalues $\lambda_i$ of $B_q$ within
\begin{equation}
    |\lambda_i-\tilde\lambda_i|\le
    \frac{2\|G\|_F\max_j|\tilde\lambda_j|+(1+\|G\|_F)\|R\|_F}{1-\|G\|_F}
    \qquad(\|G\|_F<1).
    \label{models:eq:weyl}
\end{equation}
Summing $m_qe^{-\beta\lambda_i}$ over all blocks in interval arithmetic
encloses $\log Z_{n_0}(\beta_0)$, $\log Z_{2n_0}(\beta_0)$, and
$\log Z_{2n_0}(2\beta_0)$, yielding the certified rational upper bounds
$(p_0,q_0)$. The staggered-field, dimerized, ladder, and mixed-field
Ising models in Table~\ref{models:tab:summary} come from this pipeline
directly; the single-ion anisotropy model at $D=4$ uses the half-length
lift below.

When $2n_0$ reaches the limit of block diagonalization, a direct seed at
length $2n_0$ is out of reach, since both of its purities need the
partition function of a ring of $4n_0$ cells. The spectra already in
hand at lengths $n_0$ and $2n_0$ bound $1-\Ps(n_0,\beta_0)\le s$,
$1-\Ps(n_0,2\beta_0)\le s_2$, and $1-\Pt(2n_0,\beta_0)\le t$. The
obstacle is the thermal purity at $4n_0$, which no computed spectrum
reaches; the torus identity~\eqref{sm:eq:torus} at length $2n_0$
supplies it:
\begin{proofdiagram}
    \halflengthlift
\end{proofdiagram}
Move 2 is this identity with the denominator $\Ps(2n_0,\beta_0)^2\le1$
dropped%
~\leanmark{https://sirui-lu.com/CertifiedGaps/docs/CertifiedGaps/PurityCriterion.html\#CertifiedGaps.thermalPurity_four_mul}:
\begin{equation}
    \Pt(4n_0,\beta_0)
    =\frac{\Ps(2n_0,2\beta_0)\,\Pt(2n_0,\beta_0)^2}{\Ps(2n_0,\beta_0)^2}
    \ge\bigl(1-f(s_2)\bigr)(1-t)^2.
    \label{models:eq:half-length}
\end{equation}
For the spin-$1$ chain with single-ion anisotropy $D=4$ at
$\beta_0=29/10$, the $L=6$ and $L=12$ spectra ($3^{12}=531{,}441$
states) supply $s$, $s_2$, and $t$, and the lift gives the seed at
$n_0=12$ without the $24$-site ring ($3^{24}\simeq2.8\times10^{11}$
states). It sharpens the uniform gap bound at $D=4$ from $0.6946$ (the
Lean chain from the direct seed $(6,9/4)$%
~\leanmark{https://sirui-lu.com/CertifiedGaps/docs/Leaderboard/Records/SpinOneLargeD.html\#Leaderboard.Records.SpinOneLargeD.D_4})
to $0.8994$ for all even
$L\ge12$; the short rings $L=6,8,10$ are covered directly by their
verified spectral gaps.

\begin{table*}[t]
    \caption{Seed scales $(n_0,\beta_0)$, certified starting bounds
        $p_0\ge1-\Ps(n_0,\beta_0)$ and $q_0\ge1-\Pt(2n_0,\beta_0)$, and
        resulting uniform gap bounds for all eight models in
        Table~\ref{tab:models} ($n_0$ counts sites, two-spin dimer cells, or
        two-spin rungs as in Table~\ref{tab:models}). The pairs $(p_0,q_0)$ are
        four-digit upper bounds on the enclosures of the model certificates and
        enter the Lean theorems as hypotheses; the last row is the lifted seed
        $(k,r)=(28,17/8)$ of Table~\ref{cert:tab:lifted}. The printed gap is the
        interval bound of the certificate; each linked Lean theorem proves the
        rational-chain bound~\eqref{models:eq:chain-cert} from its seed pair,
        which agrees in the quoted digits except for the single-ion row, whose
        Lean theorem starts from the direct seed $(6,9/4)$ and proves $0.6946$.
        The mixed-field Ising seed is proved unconditionally in the Lean kernel.}
    \label{models:tab:summary}
    \begin{tabular}{@{}llcccc@{}}
        \toprule
        Model                           & Parameter           & Seed $(n_0,\beta_0)$ & $p_0\ge1-\Ps(n_0,\beta_0)$ & $q_0\ge1-\Pt(2n_0,\beta_0)$ & Certified gap $\ge$                                                                                                                                               \\
        \midrule
        Staggered field, $S=\tfrac12$   & $h_{\mathrm{st}}=1$ & $(8,21/4)$           & $6.937\times10^{-4}$       & $7.327\times10^{-4}$        & $1.1489$~\leanmark{https://sirui-lu.com/CertifiedGaps/docs/CertifiedGaps/Models/StaggeredHeisenbergPhysical.html\#CertifiedGaps.StaggeredHeisenberg.physical_gap} \\
        Dimerized, $S=\tfrac12$         & $\delta=1/2$        & $(4,6)$              & $1.133\times10^{-2}$       & $1.224\times10^{-2}$        & $0.5347$~\leanmark{https://sirui-lu.com/CertifiedGaps/docs/CertifiedGaps/Models/DimerizedHeisenbergPhysical.html\#CertifiedGaps.DimerizedHeisenberg.physical_gap} \\
        Dimerized, $S=1$                & $\delta=9/10$       & $(2,7/2)$            & $4.972\times10^{-2}$       & $3.206\times10^{-2}$        & $0.5288$~\leanmark{https://sirui-lu.com/CertifiedGaps/docs/CertifiedGaps/Models/SpinOneDimerized.html\#CertifiedGaps.SpinOneDimerized.gap}                        \\
        Single-ion anisotropy, $S=1$    & $D=4$               & $(12,29/10)$         & $3.633\times10^{-3}$       & $3.141\times10^{-2}$        & $0.8994$~\leanmark{https://sirui-lu.com/CertifiedGaps/docs/CertifiedGaps/Models/SpinOneLargeD.html\#CertifiedGaps.SpinOneLargeD.gap}                              \\
        Heisenberg ladder, $S=\tfrac12$ & $J_\perp=4$         & $(4,2)$              & $3.509\times10^{-2}$       & $2.463\times10^{-2}$        & $1.1024$~\leanmark{https://sirui-lu.com/CertifiedGaps/docs/CertifiedGaps/Models/HeisenbergLadderPhysical.html\#CertifiedGaps.HeisenbergLadder.physical_gap}       \\
        Ising ladder                    & $g=3$               & $(4,3/2)$            & $1.125\times10^{-1}$       & $5.162\times10^{-2}$        & $0.6368$~\leanmark{https://sirui-lu.com/CertifiedGaps/docs/CertifiedGaps/Models/TwoLegIsing.html\#CertifiedGaps.TwoLegIsing.gap}                                  \\
        Mixed-field Ising               & $(g,h)=(1/4,1/20)$  & $(6,97/4)$           & $9.064\times10^{-6}$       & $3.141\times10^{-5}$        & $0.3922$~\leanmark{https://sirui-lu.com/CertifiedGaps/docs/Leaderboard/Records/IsingLongitudinal.html\#Leaderboard.Records.IsingLongitudinal.g_1_4_h_1_20}        \\
        Heisenberg, $S=1$               & isotropic           & $(60,833/32)$        & $3.054\times10^{-2}$       & $4.796\times10^{-2}$        & $0.0729$~\leanmark{https://sirui-lu.com/CertifiedGaps/docs/Leaderboard/Records/SpinOneHeisenberg.html\#Leaderboard.Records.SpinOneHeisenberg.temperature_lifting} \\
        \bottomrule
    \end{tabular}
\end{table*}

Finally, to verify the $J$ doubling steps~\eqref{sm:eq:update} from
$(p_0,q_0)$ and the resulting gap bound $\Delta_L>\tfrac{\log r}{\beta_0}$ in
Lean, whose kernel decides rational inequalities but has no certified
logarithm, we express the
interval bounds~\eqref{cert:eq:interval-gap} and the quadratic tail
of Sec.~\ref{sm:sec:gap} in terms of a rational base factor $r>1$ and
tail threshold $\varepsilon\le1/100$:
\begin{equation}
    r^{2^j}z_j<1-z_j\quad(0\le j<J),\qquad
    p_J,q_J\le\varepsilon\le\frac1{100},\qquad
    5\varepsilon r^{2^J}\le1,
    \label{models:eq:chain-cert}
\end{equation}
where $z_j=1-(1-p_j)^2(1-q_j)$ along the upward-rounded rational
sequence $(p_j,q_j)$. The first inequality is equivalent to
$\tfrac{\log[(1-z_j)/z_j]}{2^j\beta_0}>\tfrac{\log r}{\beta_0}$ on $[2^jn_0,2^{j+1}n_0]$,
and the remaining two initialize the quadratic tail
inequality~\eqref{sm:eq:scalar} for all $j\ge J$. The Lean kernel checks these rational
inequalities directly, and the general theorem%
~\leanmark{https://sirui-lu.com/CertifiedGaps/docs/CertifiedGaps/Ring/Gap.html\#CertifiedGaps.RPRing.gapGT_of_chainCertificate}
concludes that any reflection-positive ring satisfying
$1-\Ps(n_0,\beta_0)\le p_0$ and $1-\Pt(2n_0,\beta_0)\le q_0$ has a unique
ground state and $\Delta_L>\tfrac{\log r}{\beta_0}$ for every even $L\ge n_0$.
All eight seeds of Table~\ref{models:tab:summary} (labeled by model) lie
inside the admissible region of the asymmetric recurrence and below the
symmetric threshold $\varepsilon_*\simeq0.1273$:
\begin{proofdiagram}
    \seedscatter
\end{proofdiagram}
The boundary is the update~\eqref{sm:eq:update} iterated in floating
point until $(p_j,q_j)$ enters $[0,1/100]^2$; it admits $p_0$ up to
$0.217$ when $q_0\to0$ and $q_0$ up to $0.200$ when $p_0\to0$, so the
recurrence tolerates $1-\Ps$ and $1-\Pt$ almost independently.

\subsection{Unconditional seed proof for the mixed-field Ising chain}
\label{models:sec:ising}

Section~\ref{models:sec:spectra} supplies the seed bounds $(p_0,q_0)$ to
the Lean theorem as hypotheses. To eliminate that hypothesis and prove an
unconditional gap theorem in Lean, the kernel must bound
$Z_{n_0}(\beta_0)$, $Z_{2n_0}(\beta_0)$, and $Z_{2n_0}(2\beta_0)$ from
above and below directly. A full $2^{12}\times2^{12}$ eigensolve inside
a proof-assistant kernel is prohibitive, but at the mixed-field Ising
point $(g,h)=(1/4,1/20)$,
\begin{equation}
    H_L=-\sum_{j=1}^L\sigma_j^z\sigma_{j+1}^z
    -\frac14\sum_{j=1}^L\sigma_j^x
    -\frac1{20}\sum_{j=1}^L\sigma_j^z,
    \label{models:eq:ising}
\end{equation}
we do not need the full spectrum. On a ring of length $L\in\{6,12\}$,
the longitudinal field splits the two classical ferromagnetic states
$\ket{\uparrow\cdots\uparrow}$ and $\ket{\downarrow\cdots\downarrow}$ by
$2hL=L/10$, whereas creating a pair of domain walls costs energy $4$.
Because the transverse field $-g\sum_{j=1}^L\sigma_j^x$ has operator
norm $gL=L/4$, which at $L\le12$ exceeds the ferromagnetic splitting
$2hL$ but remains below the domain-wall cost $4$, bounding the
transverse field by its norm separates all $2^L-2$ domain-wall states
from the two ferromagnetic levels, while two integer trial vectors
locate the two lowest eigenvalues to within $\delta=10^{-10}$:
\begin{proofdiagram}
    \isingsandwich
\end{proofdiagram}

\emph{Two low-lying eigenvalues from integer trial residuals.}
For any Hermitian matrix $H$ with eigenvalues $\lambda_i$, any nonzero
trial vector $x$, and any real shift $\theta$, expanding
$\|(H-\theta\smI)x\|^2$ in the eigenbasis shows that at least one
eigenvalue $\lambda_i$ satisfies
$(\lambda_i-\theta)^2\|x\|^2\le\|(H-\theta\smI)x\|^2$. Because $H_L$ has
rational entries in the $\sigma^z$ basis, this residual inequality
reduces to an exact integer dot-product comparison. On each ring
$L\in\{6,12\}$, we supply two integer trial vectors
$x_0,x_1\in\mathbb Z^{2^L}$ scaled by $2^{52}$ and rational candidate
energies $\theta_0<\theta_1$. The Lean kernel checks in exact integer
arithmetic that%
~\leanmark{https://sirui-lu.com/CertifiedGaps/docs/SeedKernel/Ising/Seed.html\#SeedKernel.Ising.exists_near_of_check}
\begin{equation}
    \|(H_L-\theta_a\smI)x_a\|^2\le\delta^2\|x_a\|^2
    \qquad(a=0,1,\quad\delta=10^{-10}),
    \label{models:eq:residual}
\end{equation}
which places two distinct eigenvalues of $H_L$ in the disjoint intervals
$[\theta_0-\delta,\theta_0+\delta]$ and
$[\theta_1-\delta,\theta_1+\delta]$ ($\theta_0\simeq-6.3918$ and
$\theta_1\simeq-5.7966$ at $L=6$; $\theta_0\simeq-12.7835$ and
$\theta_1\simeq-11.5932$ at $L=12$).

\emph{Lower bound on the remaining $2^L-2$ eigenvalues by configuration counting.}
To verify that the two eigenvalues in $[\theta_0\pm\delta]$ and
$[\theta_1\pm\delta]$ are the two lowest levels of $H_L$ and to bound
all remaining $2^L-2$ eigenvalues from below, we split $H_L$ in the
$\sigma^z$ basis $\ket{\sigma}$ ($\sigma\in\{\pm1\}^L$) into the
diagonal Ising energy $E_{\mathrm{cl}}(\sigma)$ and the off-diagonal
transverse field $-g\sum_{j=1}^L\sigma_j^x$. Since each $\sigma_j^x$ has
operator norm one, the transverse field contributes at most $gL\|x\|^2$
to $\langle x,H_Lx\rangle$, giving the quadratic-form lower bound
\begin{equation}
    \langle x,H_Lx\rangle
    \ge\sum_{\sigma\in\{\pm1\}^L}\bigl(E_{\mathrm{cl}}(\sigma)-gL\bigr)|x(\sigma)|^2.
    \label{models:eq:gershgorin}
\end{equation}
For any threshold $\rho$, let
$B=\{\sigma:E_{\mathrm{cl}}(\sigma)-gL<\rho\}$. By the min--max
principle, at most $|B|$ eigenvalues of $H_L$ lie below $\rho$, because
Eq.~\eqref{models:eq:gershgorin} gives $\langle x,H_Lx\rangle\ge\rho\|x\|^2$
on the codimension-$|B|$ subspace of vectors vanishing on $B$. On a
ring, the two uniform configurations \spinrow{1,1,1,1,1,1} and
\spinrow{-1,-1,-1,-1,-1,-1} have classical energies $-L\mp hL$; every
other configuration, such as \spinrow{1,1,-1,-1,1,1}, has at least two
domain walls ($\sum_j\sigma_j\sigma_{j+1}\le L-4$) and at least one spin
against the field ($\sum_j\sigma_j\le L-2$), so its shifted classical
energy is at least
\begin{equation}
    \rho=-(L-4)-h(L-2)-gL,
    \label{models:eq:ising-rho}
\end{equation}
which equals $-3.7>\theta_1+\delta$ at $L=6$ and $-11.5>\theta_1+\delta$
at $L=12$. The kernel checks $|B|=2$ at this $\rho$ over all $2^L$
configurations, so, with two eigenvalues already located below
$\theta_1+\delta<\rho$, the remaining $2^L-2$ satisfy $\lambda_i\ge\rho$.

\emph{Partition-function bounds and unconditional gap theorem.}
Combining the two eigenvalue enclosures with the lower bound
$\lambda_i\ge\rho$ ($i\ge2$) sandwiches $Z_L(\beta)$ from both sides:
dropping all excited states bounds $Z_L(\beta)$ from below by the ground
state alone, $Z_L(\beta)\ge e^{-\beta(\theta_0+\delta)}$, while bounding
the two lowest eigenvalues from below by $\theta_0-\delta$ and
$\theta_1-\delta$ and all remaining $2^L-2$ eigenvalues by $\rho$ gives
the upper bound%
~\leanmark{https://sirui-lu.com/CertifiedGaps/docs/SeedKernel/Ising/Seed.html\#SeedKernel.Ising.partition_upper}
\begin{equation}
    Z_L(\beta)\le e^{-\beta(\theta_0-\delta)}
    \left[1+e^{-\beta(\theta_1-\theta_0)}
        +2^Le^{-\beta(\rho-\theta_0+\delta)}\right].
    \label{models:eq:ising-z}
\end{equation}
The leading factor $e^{-\beta\theta_0}$ cancels between numerator and
denominator in both $\Ps(6,\beta_0)=Z_{12}(\beta_0)/Z_6(\beta_0)^2$ and
$\Pt(12,\beta_0)=Z_{12}(2\beta_0)/Z_{12}(\beta_0)^2$ because
$2\theta_0(6)\simeq\theta_0(12)$ to within $10^{-7}$. Using rational
bounds on $e$ in Lean, evaluating Eq.~\eqref{models:eq:ising-z} at
$\beta_0=97/4$ proves the starting bounds at $(n_0,\beta_0)=(6,97/4)$%
~\leanmark{https://sirui-lu.com/CertifiedGaps/docs/SeedKernel/Ising/Seed.html\#SeedKernel.Ising.first_point_seed},
\begin{equation}
    p_0=\frac{9064}{10^9}\ge1-\Ps(6,97/4),\qquad
    q_0=\frac{3141}{10^8}\ge1-\Pt(12,97/4).
    \label{models:eq:ising-seed}
\end{equation}
At $j=0$, these bounds give $z_0=1-(1-p_0)^2(1-q_0)\simeq4.954\times10^{-5}$,
so $(1-z_0)/z_0>20{,}184$; substituting Eq.~\eqref{models:eq:ising-seed}
into the rational doubling inequalities~\eqref{models:eq:chain-cert}
over $J=14$ steps with base factor $r=13531935/10^3<(1-z_0)/z_0$ proves
a unique ground state and%
~\leanmark{https://sirui-lu.com/CertifiedGaps/docs/Leaderboard/Records/IsingLongitudinal.html\#Leaderboard.Records.IsingLongitudinal.g_1_4_h_1_20}
\begin{equation}
    \boxed{
        \Delta_L>\frac{\log(13531.935)}{97/4}>0.3922
        \qquad\text{for every even }L\ge6
    }
    \label{models:eq:ising-gap}
\end{equation}
unconditionally in Lean.
